\section{Methods}\label{sec:methods}

\subsection{Background}\label{sec:background}

\subsubsection{E-variables, e-processes and Ville's inequality}
\label{sec:eproc}

Let $(\Omega, \Fcal, (\Fcal_t)_{t \ge 0}, \Prob)$ be a filtered probability
space and let $H_0$ be a set of distributions on it.

\begin{definition}[E-variable]\label{def:evar}
A non-negative random variable $E$ is an \emph{e-variable} for $H_0$ if
$\Exp_P[E] \le 1$ for every $P \in H_0$.
\end{definition}

\begin{definition}[E-process]\label{def:eproc}
A non-negative, adapted process $(E_t)_{t \ge 0}$ is an \emph{e-process} for
$H_0$ if $\Exp_P[E_\tau] \le 1$ for every $P \in H_0$ and every stopping time
$\tau$ (finite or not, with $E_\infty := \liminf_t E_t$).
\end{definition}

Every non-negative supermartingale starting at one is an e-process, by
optional stopping \citep{shafer2019game}; the converse fails, and the distinction matters for
composite nulls, where an e-process need not be a supermartingale under any
single null distribution \citep{ramdas2023gametheoretic}. Both of our
constructions are e-processes that are dominated pointwise by a
$P$-supermartingale for each $P \in H_0$, with the dominating process
depending on $P$. That is the standard route to an e-process for a composite
null, and it is the route both proofs below take.

The operational content of Definition~\ref{def:eproc} is Ville's inequality
\citep{ville1939etude}: for an e-process $(E_t)$ and any $\alpha \in (0,1)$,
\begin{equation}\label{eq:ville}
\Prob\bigl( \exists\, t \ge 0 : E_t \ge 1/\alpha \bigr) \;\le\; \alpha
\qquad \text{for every } P \in H_0 .
\end{equation}
The quantifier is inside the probability. This is what licenses continuous
monitoring: an analyst may inspect the evidence after every observation and
stop whenever they wish, and the probability that the process \emph{ever}
crosses the threshold under the null is still at most $\alpha$. A p-value
offers no analogue. A fixed-sample p-value controls the rejection probability
at one pre-specified sample size, and the guarantee is destroyed by looking;
the phenomenon is familiar from sequential clinical trials and from online
experimentation \citep{johari2022always}.

E-values also compose in ways p-values do not, and we use two such
composition rules. If $E^{(1)}, \dots, E^{(m)}$ are e-variables for
$H_0^{(1)}, \dots, H_0^{(m)}$ then $\min_k E^{(k)}$ is an e-variable for
$\bigcup_k H_0^{(k)}$, and any convex combination $\sum_k w_k E^{(k)}$ with
$w_k \ge 0$, $\sum_k w_k = 1$ is an e-variable for $\bigcap_k H_0^{(k)}$,
under \emph{arbitrary} dependence between the components
\citep{vovk2021evalues}. The first rule is what makes the adjacency
certificate possible; the second is what makes the multiplicity discussion of
Section~\ref{sec:multiplicity} non-trivial.

Recent work has developed e-processes for a wide range of settings---bounded
means \citep{waudbysmith2024estimating}, linear models
\citep{lindon2026anytime}, stratified counts \citep{turner2023safe}---and has
characterised admissible constructions \citep{ramdas2023gametheoretic}. Two
general recipes recur below. The first is the \emph{method of mixtures}: under
a group-invariant null, integrating the likelihood ratio against the right-Haar
measure on the nuisance group yields an exact test martingale
\citep{berger1998bayes,hendriksen2021optional,grunwald2024safe}.
Section~\ref{sec:primitive} instantiates it for partial correlation. The second
is \emph{universal inference} \citep{wasserman2020universal}: for a null model
$\{D_\theta : \theta \in \Theta_0\}$ and any alternative fitted on past data,
the ratio of a sequentially-fitted alternative likelihood to the maximised null
likelihood is an e-process, with no regularity conditions on $\Theta_0$ beyond
the existence of the supremum. Section~\ref{sec:direction} instantiates it for
causal direction. Universal inference is known to be conservative
\citep{tse2022note}, and Section~\ref{sec:exp-orientation} measures how
conservative it is here.

\subsubsection{Causal discovery setting}\label{sec:setting}

Let $\Gstar$ be a directed acyclic graph (DAG) with vertex set
$V = \{1,\dots,d\}$ identified with random variables $X_1,\dots,X_d$, and let
$P$ be the joint law of $X = (X_1,\dots,X_d)$. We write $\pa(i)$ for the
parents of $i$ in $\Gstar$ and $\de(i)$ for its descendants. Throughout,
observations $o_1, o_2, \dots$ are drawn i.i.d.\ from $P$ and arrive in
batches; $\Fcal_t$ is the $\sigma$-field generated by the first $t$
observations.

\begin{assumption}[Markov condition]\label{ass:markov}
$P$ is Markov to $\Gstar$: whenever $S$ d-separates $i$ from $j$ in $\Gstar$,
$X_i \indep X_j \mid X_S$ under $P$.
\end{assumption}

Assumption~\ref{ass:markov} is the direction of the correspondence that
follows from the causal structure itself: it says that the graph's separations
are reflected in the distribution. Its converse---faithfulness, that every
conditional independence in $P$ corresponds to a d-separation in
$\Gstar$---does not follow from the causal structure, and is the assumption we
do not make for validity.

Write $\Scal_k(i,j) = \{ S \subseteq V \setminus \{i,j\} : |S| \le k \}$ for
the family of admissible conditioning sets. This family is determined by $d$
and $k$ and is fixed before any data are seen; nothing about it adapts to the
sample. We say the pair $(i,j)$ is \emph{separable within $\Scal_k$} under $P$
if $X_i \indep X_j \mid X_S$ for some $S \in \Scal_k(i,j)$.

\begin{assumption}[Separator size]\label{ass:sep}
Every non-adjacent pair of $\Gstar$ has a d-separating set of size at most
$k$.
\end{assumption}

Assumption~\ref{ass:sep} is structural rather than distributional: it
constrains $\Gstar$, not $P$, and it can be guaranteed by choosing $k$ large
enough. It holds whenever $k$ is at least the maximum in-degree of $\Gstar$,
since for non-adjacent $i,j$ either $j \notin \de(i)$, in which case $\pa(i)$
d-separates them, or $i \notin \de(j)$, in which case $\pa(j)$ does. Taking
$k = d-2$ makes it vacuous at the cost of an exponential family $\Scal_k$; the
trade-off is computational, not inferential, and we return to it in
Section~\ref{sec:complexity}. Section~\ref{sec:exp-premise} measures how often
Assumption~\ref{ass:sep} holds at the $k$ we use, and compares it with how
often $\delta$-strong faithfulness holds on the same instances.

\subsubsection{Faithfulness and its quantitative strengthenings}
\label{sec:faithfulness}

Since the argument of this paper is about what faithfulness is doing, it is
worth being precise about the versions in play.

\emph{Faithfulness} holds if every conditional independence in $P$ is implied
by d-separation in $\Gstar$; \citet{meek1995strong} showed that it fails only
on a Lebesgue-null set of parameters for a fixed graph, which is the result
usually cited in its defence and which \citet{uhler2013geometry} showed to be
misleading in finite samples. \emph{Adjacency-faithfulness}
\citep{ramsey2006adjacency} is the weaker requirement that adjacent variables
are dependent given every conditioning set. \emph{$\delta$-strong
faithfulness} \citep{uhler2013geometry,kalisch2007estimating} is the
quantitative strengthening that every partial correlation not forced to zero
by d-separation has absolute value at least $\delta$; it is what turns
consistency into a finite-sample statement, and it is the premise of the
comparators in Section~\ref{sec:exp-main}.

The asymmetry we exploit is this. Faithfulness in any of its forms is an
assumption about which \emph{dependences} exist. A procedure that asserts an
edge by rejecting an independence null needs it only for \emph{power}: without
it, a genuinely adjacent pair whose dependence is cancelled will fail to be
certified, and the pair stays undecided. A procedure that asserts a
\emph{non}-edge by failing to reject needs it for \emph{validity}: without it,
a cancelled dependence produces a confidently deleted edge. The same
assumption sits on opposite sides of the guarantee depending on which
direction of the inference the algorithm runs. \citet{zhang2008detection} and
\citet{ramsey2006adjacency} respond to this by detecting or tolerating
unfaithfulness within a delete-on-non-rejection architecture; we respond by
changing the architecture.

\subsection{The certificate-only principle}\label{sec:principle}

We state the rule formally and record what it forces.

\begin{definition}[Certificate]\label{def:cert}
Let $\phi$ be a feature of $\Gstar$---a proposition such as ``$i$ and $j$ are
adjacent'' or ``$\Gstar$ contains the arrowhead $i \to j$''. A
\emph{certifying null} for $\phi$ is a hypothesis $H_\phi$ with the property
that $P \in H_\phi$ whenever $\phi$ is false. A \emph{certificate for $\phi$
at level $\alpha$} is the event that an e-process for $H_\phi$ exceeds
$1/\alpha$.
\end{definition}

The implication in Definition~\ref{def:cert} runs one way only, and that is
the whole point. $H_\phi$ must be true whenever $\phi$ is false; it may also
be true in some cases where $\phi$ is true, and then the certificate simply
never fires. This asymmetry is what buys validity without faithfulness, and
what costs power.

\begin{proposition}[Soundness of a single certificate]\label{prop:single}
If $(E_t)$ is an e-process for a certifying null $H_\phi$, then
$\Prob(\exists t: \text{$\phi$ is certified at level $\alpha$ and $\phi$ is
false}) \le \alpha$.
\end{proposition}

\begin{proof}
If $\phi$ is false then $P \in H_\phi$ by Definition~\ref{def:cert}, so
$(E_t)$ is an e-process under $P$ and~\eqref{eq:ville} bounds the probability
that it ever exceeds $1/\alpha$ by $\alpha$. If $\phi$ is true the event in
question is empty.
\end{proof}

Proposition~\ref{prop:single} is elementary; we state it to make explicit that
the burden of the design is entirely in constructing $H_\phi$, not in the
inference. Three structural consequences follow, and they shape everything
below.

\paragraph{The algorithm cannot delete.} There is no certifying null for the
proposition ``$i$ and $j$ are non-adjacent''. Such a null would have to be
true whenever the pair \emph{is} adjacent, that is, it would have to be a
dependence hypothesis rejected by evidence of independence; but evidence of
independence is exactly what a hypothesis test does not provide. Equivalence
testing supplies the nearest available substitute---one can reject ``the
partial correlation exceeds $\delta$ in absolute value''---but the resulting
statement is ``the association is smaller than $\delta$'', which is a claim
about effect size, not about absence. We return to this in
Section~\ref{sec:negligible}. In the absence of such a null, an algorithm
obeying the rule must start from the empty graph.

\paragraph{The output is three-valued.} Each pair is in exactly one of three
states: \emph{certified adjacent} (with or without a certified orientation),
\emph{undecided}, or---if the extension of Section~\ref{sec:negligible} is
enabled---\emph{certified negligible at margin $\delta$}. The absence of an
edge in the output graph is not an assertion.

\paragraph{Power is descriptive, not guaranteed.} Nothing in
Proposition~\ref{prop:single} promises that a true edge is ever certified.
Whether it is depends on adjacency-faithfulness, on the sample size, and on
the primitive's power. This is a real limitation and we treat it as one:
Section~\ref{sec:experiments} reports certification rates alongside error
rates throughout, and reports the fraction of pairs left undecided rather than
suppressing it.

\subsection{The adjacency certificate}\label{sec:adjacency}

\subsubsection{The separability null}\label{sec:sepnull}

Fix a pair $(i,j)$ and consider
\begin{equation}\label{eq:sepnull}
H^{\mathrm{sep}}_{ij} \;:\quad
\exists\, S \in \Scal_k(i,j) \ \text{ such that } \ X_i \indep X_j \mid X_S .
\end{equation}

\begin{lemma}\label{lem:certifying}
Under Assumptions~\ref{ass:markov} and~\ref{ass:sep},
$H^{\mathrm{sep}}_{ij}$ is a certifying null for the adjacency of $i$ and $j$.
\end{lemma}

\begin{proof}
Suppose $i$ and $j$ are non-adjacent in $\Gstar$. By
Assumption~\ref{ass:sep} there is a set $S$ with $|S| \le k$ that d-separates
them, and $S \subseteq V \setminus \{i,j\}$, so $S \in \Scal_k(i,j)$. By
Assumption~\ref{ass:markov}, $X_i \indep X_j \mid X_S$ under $P$. Hence
$P \in H^{\mathrm{sep}}_{ij}$.
\end{proof}

Note what the lemma does \emph{not} require. It says nothing about adjacent
pairs. If $i$ and $j$ are adjacent, $P$ may or may not lie in
$H^{\mathrm{sep}}_{ij}$: an adjacent pair whose dependence is exactly
cancelled by a path through some $S$ satisfies the null too, and then the
certificate cannot fire. That is the failure of adjacency-faithfulness, and it
costs power, not validity.

\subsubsection{The minimum construction}\label{sec:minconstruction}

Let $(E^{(S)}_t)_{t \ge 0}$ be an e-process for the point null
$X_i \indep X_j \mid X_S$, for each $S \in \Scal_k(i,j)$, all built from the
same data stream. Define
\begin{equation}\label{eq:adjacency}
A_t(i,j) \;=\; \min_{S \in \Scal_k(i,j)} E^{(S)}_t(i,j) .
\end{equation}

\begin{proposition}\label{prop:adjacency}
$(A_t)_{t \ge 0}$ is an e-process for $H^{\mathrm{sep}}_{ij}$.
\end{proposition}

\begin{proof}
Let $P \in H^{\mathrm{sep}}_{ij}$. By~\eqref{eq:sepnull} there exists
$S^\star \in \Scal_k(i,j)$ with $X_i \indep X_j \mid X_{S^\star}$ under $P$,
so $P$ satisfies the point null for which $(E^{(S^\star)}_t)$ is an e-process.
By~\eqref{eq:adjacency}, $A_t \le E^{(S^\star)}_t$ pointwise for every $t$,
hence $A_\tau \le E^{(S^\star)}_\tau$ for every stopping time $\tau$, and
non-negativity gives
$\Exp_P[A_\tau] \le \Exp_P[E^{(S^\star)}_\tau] \le 1$.
\end{proof}

Combining Lemma~\ref{lem:certifying}, Proposition~\ref{prop:adjacency} and
Proposition~\ref{prop:single} gives the adjacency guarantee: the probability
that $A_t(i,j)$ ever exceeds $1/\alpha$ for a non-adjacent pair is at most
$\alpha$.

That a minimum of e-values is an e-value for a union null is standard
\citep{vovk2021evalues}, and the proof above is the standard one. What is
specific here is the identification of adjacency as the complement of such a
union, and the two consequences that follow.

\begin{remark}[No faithfulness]\label{rem:faithfulness}
The proof uses only that~\eqref{eq:sepnull} is true for non-adjacent pairs,
which follows from Assumptions~\ref{ass:markov} and~\ref{ass:sep}.
Faithfulness enters only through power: adjacency-faithfulness
\citep{ramsey2006adjacency} is what makes every $E^{(S)}$ grow for a genuinely
adjacent pair, and hence what makes the minimum grow. An instance where it
fails leaves pairs undecided rather than producing wrong edges. Stated
exactly, validity requires three things: the Markov condition,
Assumption~\ref{ass:sep}, and a valid per-set e-process. The third is not
free. Our default primitive (Section~\ref{sec:primitive}) is a linear-Gaussian
safe test, so that model is assumed \emph{there}; substituting a nonparametric
sequential test \citep{he2026sequential} removes it at a cost in power.
Proposition~\ref{prop:adjacency} is indifferent to which primitive is used,
and Section~\ref{sec:exp-primitive} measures the primitive's calibration under
six noise laws, including three outside its model.
\end{remark}

\begin{remark}[Multiplicity over pairs]\label{rem:multiplicity-pairs}
The construction tests one hypothesis per pair, not one per (pair,
conditioning set) triple. A union bound over the adjacency family therefore
runs over $\binom{d}{2}$ hypotheses rather than
$\binom{d}{2} \cdot |\Scal_k|$. At $d = 6$ and $k = 2$ we have
$|\Scal_k| = \binom{4}{0} + \binom{4}{1} + \binom{4}{2} = 11$, so the
per-hypothesis level is $11$ times larger than a triple-indexed budget would
allow. The saving grows with $k$: at $d = 10$, $k = 3$ it is $93$. This is not
a free lunch---the minimum in~\eqref{eq:adjacency} is itself conservative,
since it discards all but the least favourable conditioning set---but the
conservatism is of a different kind from a multiplicity correction, and it is
the kind that vanishes as the least favourable set becomes uninformative.
\end{remark}

\subsubsection{The per-set primitive}\label{sec:primitive}

Proposition~\ref{prop:adjacency} needs an e-process for each point null
$X_i \indep X_j \mid X_S$. For linear-Gaussian data we use an exact one,
obtained as a Bayes factor with right-Haar priors on the nuisance parameters.
Write the hypothesis in regression form,
\begin{equation}\label{eq:regform}
X_i \;=\; X_S^\top \alpha \;+\; \beta X_j \;+\; \varepsilon,
\qquad \varepsilon \sim N(0, \sigma^2),
\end{equation}
so that the null is $\beta = 0$ with nuisance parameters $(\alpha, \sigma)$
shared between the hypotheses. Place the right-Haar prior
$\pi(\alpha, \sigma) \propto \sigma^{-1}$ on the nuisance block under both
hypotheses, and $\beta \mid \sigma \sim N(0, \sigma^2 g)$ on the tested
coefficient under the alternative, in the manner of Zellner's $g$-prior
\citep{zellner1986assessing}. The right-Haar choice on the nuisance block is
the invariance-based prior of \citet{jeffreys1946invariant} specialised to
this group. Integrating both marginal likelihoods in
closed form (Appendix~\ref{app:safelinear}) gives, after $n$ observations,
\begin{equation}\label{eq:safelinear}
E_n \;=\; (1+G)^{-1/2}\,\bigl(1 - \kappa\,\hat r_n^2\bigr)^{-(n-p)/2},
\qquad
G \;=\; g\,S_{yy}, \quad \kappa \;=\; \frac{G}{1+G},
\end{equation}
where $p = |S| + 1$; $S_{xx}$, $S_{xy}$, $S_{yy}$ are the cross-products of
the $X_S$-residualised versions of $X_j$ and $X_i$; and
$\hat r_n^2 = S_{xy}^2 / (S_{xx} S_{yy})$ is the squared sample partial
correlation.

\begin{proposition}\label{prop:safelinear}
Under the linear-Gaussian model~\eqref{eq:regform} with $\beta = 0$, the
process $(E_n)$ of~\eqref{eq:safelinear} is a test martingale with respect to
the filtration generated by the observations, and hence an e-process.
\end{proposition}

The proof (Appendix~\ref{app:safelinear}) is the standard right-Haar argument
\citep{berger1998bayes,hendriksen2021optional}: the null model is invariant
under the group of location shifts in the $X_S$ directions and scale changes,
the right-Haar prior is relatively invariant under that group, and the
resulting Bayes factor is a ratio of marginal likelihoods of the maximal
invariant, whose null distribution is free of the nuisance parameters. Bayes
factors of this form are exactly the ones for which optional stopping is
valid, which is why not every Bayes factor would do here.

Two properties of~\eqref{eq:safelinear} drive the implementation.

\paragraph{Sufficiency.} $E_n$ depends on the data only through Schur
complements of the running Gram matrix
$\sum_{s \le n} (o_s - \bar o)(o_s - \bar o)^\top$. A single $O(d^2)$
sufficient statistic therefore serves the \emph{entire} hypothesis family
$\{(i,j,S)\}$ simultaneously: no per-hypothesis state is stored, and the cost
of an update is $O(d^2)$ regardless of how many hypotheses are live. This is
what makes it feasible to keep every pair's minimum
in~\eqref{eq:adjacency} current after every batch.

\paragraph{The prior scale must be fixed in advance.} $g$ enters
through $G = g S_{yy}$, and $S_{yy}$ grows with $n$. A prior scale chosen to
adapt to the running sample size---for instance by rescaling $g$ so that $G$
stays constant---destroys the martingale property, because the ``prior'' would
then depend on the data it is being used to weigh. We fix $g$ from a warm-up
prefix of the stream that is excluded from every e-process, using it only to
set standardisation constants. The warm-up length appears in the experiments
as a tuning parameter and nowhere in the guarantee.

\begin{remark}[Calibration was verified, not assumed]\label{rem:calibration}
Proposition~\ref{prop:safelinear} is a theorem, but the implementation of a
theorem can be wrong, and a miscalibrated primitive would invalidate every
downstream claim while leaving all the proofs intact. We therefore verified
calibration by simulation, reproducible by
\texttt{experiments/exp1\_type1\_calibration.py --reps 3000 --nmax 1500}:
monitoring~\eqref{eq:safelinear} at a grid of sample sizes over $3000$
replications up to $n = 1500$ under a true null gave realised crossing rates
of $0.108$, $0.051$, $0.025$ and $0.006$ against nominal
$\alpha \in \{0.20, 0.10, 0.05, 0.01\}$---in every case comfortably inside the
nominal level, as Proposition~\ref{prop:safelinear} requires.
Section~\ref{sec:exp-primitive} repeats the exercise under six noise laws.
\end{remark}

\subsubsection{Confidence sequence by inversion}\label{sec:cs}

Because~\eqref{eq:safelinear} is available in closed form for every value of
the tested coefficient, it can be inverted to give a confidence sequence---a
sequence of intervals covering $\beta$ simultaneously at all times with
probability $1-\alpha$ \citep{howard2021timeuniform}. Writing $\hat\beta = S_{xy}/S_{yy}$ for the
least-squares estimate of the coefficient of $X_j$ in the residualised
regression, the level-$(1-\alpha)$ confidence sequence is
\begin{equation}\label{eq:cs}
\hat\beta \;\pm\; \frac{1}{S_{yy}}
  \sqrt{\frac{\tau\,\bigl(S_{xx} S_{yy} - S_{xy}^2\bigr)}{1-\tau}},
\qquad
\tau \;=\; \frac{1}{\kappa}
  \left[ 1 - \left(\frac{\alpha}{\sqrt{1+G}}\right)^{2/(n-p)} \right] ,
\end{equation}
valid whenever $\tau > 0$, and the whole line otherwise. The derivation is in
Appendix~\ref{app:cs}. We report it because it is the natural summary of
evidence about an edge's strength, because it is what the equivalence-testing
extension of Section~\ref{sec:negligible} uses, and because it provides an
independent check on the implementation: \eqref{eq:cs} must agree with a
brute-force inversion of~\eqref{eq:safelinear}, and an early version of our
code did not. The discrepancy was a sign error in the $\log(1+G)$ term, which
produced intervals roughly half the correct width---a defect invisible in any
Type-I error simulation of the test itself, since the test was correct and
only the reported interval was wrong. Agreement between the closed form and a
numerical inversion is now a unit test.

\subsubsection{Beyond the linear-Gaussian primitive}\label{sec:beyondlg}

Conditional independence testing is hard in general: \citet{shah2020hardness}
show that no test can have non-trivial power against all alternatives while
controlling size over all null distributions with continuous conditional
laws. Every usable primitive therefore buys power with a model. The two
mainstream routes are the Model-X framework
\citep{candes2018panning}, which assumes the conditional law of $X_i$ given
$X_S$ is known, and parametric or kernel-based assumptions on the joint law.
Sequential conditional-independence tests to date have largely taken the
Model-X route; \citet{csillag2025prediction} note that the requirement is
typically inaccessible in a discovery setting, where the conditional laws are
precisely what is unknown. Construction~\eqref{eq:safelinear} avoids it by
assuming linearity and Gaussianity instead, which is a strong assumption but a
checkable one. Section~\ref{sec:exp-violations} reports what happens when it
fails.

\subsection{The direction certificate}\label{sec:direction}

\subsubsection{Why a direction can be a null hypothesis}\label{sec:whynull}

Within a Markov equivalence class no observational statistic distinguishes the
member DAGs, so ``the direction is $j \to i$'' is not a testable proposition
there: both directions have the same likelihood under every distribution. The
statement becomes testable exactly when a functional restriction breaks the
symmetry, and the classical such restriction is linear non-Gaussianity.

\begin{proposition}[\citealp{shimizu2006linear}, after
\citealp{darmois1953analyse} and \citealp{comon1994independent}]
\label{prop:lingam}
Let $X_i = \beta X_j + e_i$ with $X_j \indep e_i$, $\Exp[e_i] = 0$ and
$\beta \neq 0$. If at most one of $X_j, e_i$ is Gaussian, there is no
$\delta, e_j$ with $X_j = \delta X_i + e_j$ and $X_i \indep e_j$. If both are
Gaussian, such a representation always exists.
\end{proposition}

Proposition~\ref{prop:lingam} is the identifiability statement that LiNGAM
rests on, and it says two things we use in opposite ways. Under
non-Gaussianity, exactly one of the two factorisations of the pair is
compatible with the model, so the wrong one is a genuine, false, testable
hypothesis. Under Gaussianity, both are compatible, so no evidence can
separate them and any method that reports a direction anyway is reporting an
artefact of its tie-breaking rule. Our construction inherits both halves.

The certifying null of Section~\ref{sec:dirnull} is conditional on a
covariate block $Z$, and Proposition~\ref{prop:lingam} is bivariate and
unconditional; an earlier version of this paper cited the marginal
non-Gaussianity of $X_j$ and $e_i$ as the licence for the conditional
construction of Eqs.~\eqref{eq:fwd}--\eqref{eq:bwd}, and that citation does not
bridge the gap on its own. The bridge is short, and we record it as a lemma
because both Proposition~\ref{prop:direction} and
Proposition~\ref{prop:abstain} depend on it.

\begin{lemma}[Conditional identifiability]\label{lem:cond-lingam}
Suppose $(X_i, X_j)$ given $Z$ follows the $j \to i$ factorisation
$D^{j\to i}_\theta$ of~\eqref{eq:fwd}, with $u \indep Z$ and
$e \indep (X_j, Z)$, and write $\tilde e = \beta u + e$ for the population
$Z$-residual of $X_i$. If at most one of $u, e$ is Gaussian, then
$(X_i, X_j)$ given $Z$ does not also admit the reversed factorisation
$D^{i \to j}_\theta$ of~\eqref{eq:bwd} for any $\theta$. If both are
Gaussian, it admits both.
\end{lemma}

\begin{proof}
Because $e \indep (X_j, Z)$ and $u$ is $\sigma(X_j, Z)$-measurable, $u \indep e$,
and $\tilde e = \beta u + e$ is exactly the bivariate model of
Proposition~\ref{prop:lingam} with $u$ in the role of its ``$X_j$'' and $e$ in
the role of its ``$e_i$''. If $(X_i, X_j)$ given $Z$ also admitted
$D^{i \to j}_\theta$ for some $\theta = (b, \delta, \eta, \cdot)$, the same
computation applied to~\eqref{eq:bwd} gives that the $Z$-residual of $X_i$ is
$v := X_i - b^\top Z$ and the $Z$-residual of $X_j$ is $\delta v + w$, with
$v \indep w$. Population linear regression residuals on $Z$ are unique
(given $Z$ full rank with finite second moments), so $v = \tilde e$ and
$u = \delta \tilde e + w$ with $\tilde e \indep w$---precisely the
representation Proposition~\ref{prop:lingam} rules out when at most one of
$u, e$ is Gaussian, and permits when both are.
\end{proof}

Lemma~\ref{lem:cond-lingam} is what licenses~\eqref{eq:fwd}--\eqref{eq:bwd}: the
non-Gaussianity the construction needs is of the $Z$-residuals $u$ and $e$, not
of $X_j$ and $e_i$ themselves, and conditioning on $Z$ leaves the bivariate
argument intact only because $Z$ enters both factorisations linearly and
$e \indep (X_j, Z)$ rather than merely $e \indep X_j$.

DirectLiNGAM \citep{shimizu2011directlingam} and the pairwise measures of
\citet{hyvarinen2013pairwise} exploit the first half by \emph{comparing} the
two directions' scores and taking the larger. \citet{shimizu2008use} take a
step towards inference by testing non-normality-based direction hypotheses in
a structural equation modelling framework, and \citet{komatsu2010assessing}
assess LiNGAM's reliability by multiscale bootstrap. What we add is a
sequential e-process, which converts the comparison into a family-wise,
time-uniform error statement over a whole graph.

\subsubsection{The direction null and its e-process}\label{sec:dirnull}

Fix an ordered pair $(i,j)$ whose adjacency has been certified, and a
conditioning block $Z$ (Section~\ref{sec:freezing} explains how $Z$ is chosen
and why it is then frozen). The two competing factorisations of $(X_i, X_j)$
given $Z$ are
\begin{align}
D^{j \to i} &: \quad X_j = a^\top Z + u, \qquad
                     X_i = \beta X_j + \gamma^\top Z + e, \label{eq:fwd}\\
D^{i \to j} &: \quad X_i = b^\top Z + v, \qquad
                     X_j = \delta X_i + \eta^\top Z + w, \label{eq:bwd}
\end{align}
with $u \indep Z$, $e \indep (X_j, Z)$ and correspondingly for the second.
Both are conditional densities of the pair $(X_i, X_j)$ given $Z$, so the
marginal law of $Z$ is a common factor and cancels from any ratio; this is why
we may work with the conditional likelihoods and never model $Z$.

The certifying null for the arrowhead $i \to j$ is
\begin{equation}\label{eq:dirnull}
H^{\mathrm{dir}}_{j \to i} \;:\quad
\text{the pair was generated by the } j \to i \text{ factorisation}
\;=\; \{D^{j\to i}_\theta : \theta \in \Theta\},
\end{equation}
where $\theta = (a, \beta, \gamma, \text{noise parameters})$. Under
Lemma~\ref{lem:cond-lingam} and non-Gaussianity of the $Z$-residuals $u, e$,
$H^{\mathrm{dir}}_{j \to i}$ is true whenever the arrowhead $i \to j$ is
absent, so it is a certifying null in the sense of Definition~\ref{def:cert}.

We test it by \emph{sequential universal inference}
\citep{wasserman2020universal}. Let $\hat\theta_{s-1}$ be any estimator of the
alternative's parameters computed from $o_1, \dots, o_{s-1}$, and set
\begin{equation}\label{eq:direction}
\log E_t \;=\; \sum_{s \le t} \log D^{i \to j}_{\hat\theta_{s-1}}(o_s)
\;-\; \sup_{\theta \in \Theta} \sum_{s \le t} \log D^{j \to i}_{\theta}(o_s).
\end{equation}

\begin{proposition}\label{prop:direction}
$(E_t)_{t \ge 0}$ defined by~\eqref{eq:direction} is an e-process for
$H^{\mathrm{dir}}_{j \to i}$, provided the supremum in the second term is
attained (or exceeded) by the fitting procedure.
\end{proposition}

\begin{proof}
Let $P = D^{j \to i}_{\theta^\star} \in H^{\mathrm{dir}}_{j \to i}$ and define
\begin{equation*}
M_t \;=\; \prod_{s \le t}
  \frac{D^{i \to j}_{\hat\theta_{s-1}}(o_s)}{D^{j \to i}_{\theta^\star}(o_s)},
  \qquad M_0 = 1 .
\end{equation*}
Since $\hat\theta_{s-1}$ is $\Fcal_{s-1}$-measurable and
$D^{i \to j}_{\hat\theta_{s-1}}(\cdot)$ is a probability density for each
fixed value of $\hat\theta_{s-1}$,
\begin{equation*}
\Exp_P\!\left[ \frac{D^{i \to j}_{\hat\theta_{s-1}}(o_s)}
  {D^{j \to i}_{\theta^\star}(o_s)} \,\middle|\, \Fcal_{s-1} \right]
= \int D^{i \to j}_{\hat\theta_{s-1}}(o) \, do \;\le\; 1,
\end{equation*}
so $(M_t)$ is a non-negative supermartingale (a martingale when the densities
are mutually absolutely continuous) with $M_0 = 1$. The subtracted term
in~\eqref{eq:direction} is a supremum over $\Theta$ and $\theta^\star \in
\Theta$, so it is at least $\sum_{s\le t} \log D^{j\to i}_{\theta^\star}(o_s)$,
giving $E_t \le M_t$ pointwise. Optional stopping yields
$\Exp_P[E_\tau] \le \Exp_P[M_\tau] \le 1$ for every stopping time $\tau$.
\end{proof}

The argument is the standard universal-inference one, and it is worth being
clear that its strength is also its cost. It requires no regularity conditions
on $\Theta$, no asymptotics, and no correctly specified alternative---the
numerator may be fitted however one likes, since it only needs to integrate to
one. In exchange, subtracting a maximised likelihood rather than an integrated
one makes the resulting e-process conservative, typically by a factor related
to the null model's dimension \citep{tse2022note}.
Section~\ref{sec:exp-orientation} measures this: the certificate orients
$88\%$ of the edges DirectLiNGAM orients, which is the price of attaching an
error statement to the answer.

\subsubsection{Fitting the null, and the requirement that the supremum be real}
\label{sec:supremum}

We fit both factorisations with generalised Gaussian (exponential power)
noise,
\begin{equation}\label{eq:gg}
f_{\sigma,\kappa}(r) \;=\;
  \frac{\kappa}{2\sigma\,\Gamma(1/\kappa)}
  \exp\!\left\{ -\left|\frac{r}{\sigma}\right|^{\kappa} \right\},
\qquad \sigma > 0, \ \kappa > 0,
\end{equation}
a family \citep{subbotin1923law,box1973bayesian} whose shape parameter spans
Laplace ($\kappa = 1$), Gaussian ($\kappa = 2$) and, as
$\kappa \to \infty$, the uniform limit. Given $\kappa$, the coefficients are
fitted by iteratively reweighted least squares on the $|r|^\kappa$ loss and
$\sigma$ is profiled out in closed form; $\kappa$ is then chosen to maximise
the profile likelihood.

Proposition~\ref{prop:direction} makes a demand on this last step that is easy
to overlook, and we state it as a remark because we violated it.

\begin{remark}[The supremum must be attained]\label{rem:supremum}
The inequality $E_t \le M_t$ in the proof holds because the subtracted term is
at least the log-likelihood at $\theta^\star$. If the fitting procedure
maximises over a strict subset of $\Theta$ that excludes $\theta^\star$, the
subtracted term can fall \emph{below} the log-likelihood at $\theta^\star$,
$E_t$ can exceed $M_t$, and the e-process property fails---silently, since
nothing in the output signals it.

Our first implementation profiled $\kappa$ over the fixed grid
$\{0.7, 0.9, 1.2, 1.5, 2.0, 3.0, 5.0\}$, which omits $\kappa = 1$. On a
correctly specified Laplace null, the fitted ``supremum'' fell below the
likelihood at the true parameter in $32$ of $200$ replications, with a worst
shortfall of $2.397$ nats---an inflation of $E_t$ by a factor of up to $11$,
and hence a certificate issued at a true level of up to $11\alpha$. The fix is
to treat the grid as a bracketing device only and refine $\kappa$ continuously
by bounded scalar optimisation of the profile likelihood; the same check then
gives $0$ of $200$ violations and zero shortfall. We report this because the
defect is generic to universal inference with a discretised nuisance search,
because it is invisible to a Type-I error simulation run at the same grid (the
grid is wrong for the truth, not for itself), and because the correct check is
cheap: simulate under the null at a known $\theta^\star$ and verify that the
fitted supremum never falls below the likelihood there.
\end{remark}

\subsubsection{Abstention under Gaussianity}\label{sec:abstention}

\begin{proposition}\label{prop:abstain}
Suppose $(X_i, X_j)$ given $Z$ is jointly Gaussian, linear in $Z$. Then
$H^{\mathrm{dir}}_{j \to i}$ and $H^{\mathrm{dir}}_{i \to j}$ both hold, and
for every $\alpha \in (0,1)$ and every stopping time $\tau$, the arrowhead
$i \to j$ is certified at level $\alpha$ with probability at most $\alpha$,
and likewise for $j \to i$.
\end{proposition}

\begin{proof}
By the ``both Gaussian'' clause of Lemma~\ref{lem:cond-lingam}, applied with
$u, e$ jointly Gaussian, $(X_i, X_j)$ given $Z$ admits \emph{both}
factorisations~\eqref{eq:fwd} and~\eqref{eq:bwd} for some parameter values, so
$P \in H^{\mathrm{dir}}_{j \to i} \cap H^{\mathrm{dir}}_{i \to j}$. The proof
of Proposition~\ref{prop:direction} did not use non-Gaussianity anywhere: it
holds for every $P$ in the null it names, whichever that is. Applying it to
$H^{\mathrm{dir}}_{j \to i}$, $(E_t)$ built from~\eqref{eq:direction} is an
e-process for that null under joint Gaussianity, and Ville's
inequality~\eqref{eq:ville} gives $P(\exists\, t : E_t \ge 1/\alpha) \le \alpha$,
which is exactly the certification of $i \to j$. The symmetric argument, with
the two factorisations' roles exchanged, bounds the certification of
$j \to i$.
\end{proof}

The earlier version of this proposition claimed the deterministic statement
$\log E_t \le 0$ for every $t$, on the ground that the null and alternative
models in~\eqref{eq:direction} have identical maximised likelihoods under joint
Gaussianity. That step does not go through: the numerator and denominator are
each fitted over the generalised-Gaussian family of Eq.~\eqref{eq:gg}, whose
member distributions coincide with each other only at the single point
$\kappa = 2$, so their \emph{sample} maximised log-likelihoods need not agree
even when the data are exactly Gaussian---one family can fit an
above-Gaussian-tailed realisation better than the other does, purely by
chance. Simulating $300$ replications of a conditionally Gaussian pair at
$n = 400$ against the fitting procedure of Section~\ref{sec:supremum}, the
reversed factorisation's maximised log-likelihood exceeded the forward one in
$154$ of $300$ replications, by up to $3.8$ nats---i.e.\ $\log E_t$ does
occasionally exceed zero along the way. The revised proposition needs none of
this: it does not claim the process never rises, only that crossing its
threshold is no more likely than $\alpha$, which is what
Proposition~\ref{prop:direction} already gives for any null the pair
satisfies, Gaussian or not.

The procedure therefore abstains exactly where the direction is not
identified, at the same level $\alpha$ that governs every other certificate in
this paper, rather than deterministically or because a Gaussianity pre-test
told it to. This matters practically: a pre-test would itself have a Type-II
error, and a method that orients whenever the pre-test fails to reject
Gaussianity is once again asserting on non-rejection.
\citet{ding2025gaussdetect} give the underlying equivalence in a usable
form---Gaussianity of the forward model's noise is equivalent to independence
between regressor and residual in the \emph{reverse} regression---and use it
to avoid an explicit Gaussianity test in a different way.
Section~\ref{sec:exp-orientation} shows the abstention empirically: under
Gaussian noise the certificate orients $0$ of $692$ edges over $200$
replications, while DirectLiNGAM orients all $692$ and gets $356$ of them
wrong.

\subsubsection{Freezing the conditioning block}\label{sec:freezing}

The block $Z$ in~\eqref{eq:fwd}--\eqref{eq:bwd} is chosen when a pair's
adjacency is certified: it is the set of that pair's already-certified
neighbours at that moment. The direction process for the pair then begins
accumulating, and $Z$ is never changed.

For Eqs.~\eqref{eq:fwd}--\eqref{eq:bwd} to be the correct linear
non-Gaussian model for the pair, $Z$ must satisfy two conditions, not one:
it must contain every parent of both endpoints, \emph{and} it must exclude
every descendant of either endpoint. The first is the condition the
implementation targets and Section~\ref{sec:exp-main} measures; the second
is not targeted at all. Conditioning on a descendant of $i$ or $j$---a
collider on the path through that descendant, or a common effect of
both---induces exactly the kind of dependence a linear non-Gaussian model
does not have, for the same reason conditioning on a collider breaks
faithfulness elsewhere in this literature. The certified-neighbour freeze
does not rule this out: a variable certified adjacent to $i$ or $j$ before
the pair itself is certified can be a child of one of them, not only a
parent. Section~\ref{sec:exp-main}'s diagnostic therefore measures only
half of what correctness of the frozen null requires, and its
"complete"/"incomplete" label should be read as parent-completeness, not
as a certificate that the direction null is correctly specified.

Freezing is not a convenience. A conditioning block that kept changing would
change the hypothesis under test, and Ville's inequality bounds the crossing
probability for a \emph{fixed} null. There is a second, subtler issue:
$Z$ is chosen using the data, since it depends on which adjacencies have been
certified. The resulting guarantee is therefore conditional on the
$\sigma$-field $\Fcal_T$ at the (stopping) time $T$ at which the pair's
adjacency is certified. Because the direction process uses only observations
after $T$, and these are independent of $\Fcal_T$ given the model, the
supermartingale argument in Proposition~\ref{prop:direction} goes through
conditionally on $\Fcal_T$, and the unconditional bound follows by taking
expectations. We make this precise in Appendix~\ref{app:conditioning}. The
consequence for the implementation is that the direction process must discard
every observation before $T$, which costs power at pairs certified late.

\begin{remark}[An implementation constraint that was a bug first]
\label{rem:samesupport}
The observations scored by the numerator of~\eqref{eq:direction} and those
entering the denominator's maximised likelihood must be the \emph{same} index
set. A denominator ranging over additional observations subtracts their
log-density with no matching numerator term, which adds to $\log E_t$ a
quantity of order the number of extra observations. In our testing this
certified \emph{both} directions of every pair, including under Gaussian
noise, where Proposition~\ref{prop:abstain} guarantees the opposite. The
symptom is diagnostic: whenever a direction procedure certifies both
directions of a pair, the index sets are misaligned.
\end{remark}

\begin{remark}[Misspecification is the main theoretical risk]
\label{rem:misspec}
Proposition~\ref{prop:direction} needs the denominator's supremum to dominate
the likelihood at the true null parameter. When the noise lies outside the
fitted family~\eqref{eq:gg}, the supremum is taken over a model that does not
contain the truth and no such domination is guaranteed. This is the
certificate's principal theoretical gap, and unlike Remark~\ref{rem:supremum}
it cannot be closed by better optimisation.

Empirically the certificate is more robust than that suggests, and the
original version of this check was too weak to say so: $0$ of $12$ runs at a
single coefficient of $0.9$---an easy design, since a coefficient that large
makes the correct direction's evidence overwhelming regardless of the noise
law---has a Wilson upper limit of $0.247$, which cannot rule out a true
error rate five times $\alpha$.
\texttt{scripts/check\_misspecification\_robustness.py} reruns the same three
noise laws breaking each property of the family in turn---Student-$t_3$
(polynomial rather than exponential tails), shifted exponential (asymmetric)
and a two-component mixture (bimodal)---at $150$ replications per cell and a
coefficient sweep of $0.9$, $0.6$ and $0.3$, the last deliberately weakening
the signal the earlier check left untested. The reversed direction was
certified in $0$ of $150$ runs in every one of the nine (noise law,
coefficient) cells, with a worst-case Wilson upper limit of $0.025$ against
$\alpha/2 = 0.025$---an order of magnitude tighter than the original check,
and still zero at the weaker signal. But zero-count bounds remain bounds, not
proofs, the study is bivariate, and Section~\ref{sec:exp-sachs} shows what
happens on real data where the linear model itself is wrong: the certificate
is confidently and systematically inverted. We regard bounding the
misspecification gap analytically as the most important open problem the
construction raises.
\end{remark}

\subsection{CERT-CD}\label{sec:algorithm}

\subsubsection{The algorithm}\label{sec:algo-desc}

CERT-CD maintains a three-valued status for every unordered pair and a
two-valued status for every ordered pair. It streams data in batches. After
each batch it updates the running Gram matrix, recomputes $A_t(i,j)$ for every
undecided pair, and certifies adjacency when the running maximum of $A_t(i,j)$
crosses the pair's threshold. At that moment the pair's conditioning block is
frozen and its two direction processes are opened. Certified arrowheads are
added when a direction process crosses its own threshold.

Budgets are fixed before any data arrive. The total budget splits as
$\alpha = \alpha_A + \alpha_D$; $\alpha_A$ is divided evenly over the
$\binom{d}{2}$ adjacency nulls and $\alpha_D$ evenly over the $d(d-1)$
direction nulls, so the thresholds are
\begin{equation}\label{eq:thresholds}
\lambda_A = \frac{\binom{d}{2}}{\alpha_A},
\qquad
\lambda_D = \frac{d(d-1)}{\alpha_D}.
\end{equation}
Nothing about the split adapts to the data. Which certificates the algorithm
happens to evaluate---it only opens direction processes for certified
pairs---does not enter the bound, so the graph-dependent order in which pairs
are examined is free.

\begin{algorithm}[t]
\caption{CERT-CD}\label{alg:certcd}
\begin{algorithmic}[1]
\Require stream $o_1, o_2, \dots$; budget $\alpha$; order $k$; warm-up $m$
\State $\alpha_A, \alpha_D \gets$ split of $\alpha$;
       $\lambda_A, \lambda_D$ as in~\eqref{eq:thresholds}
\State consume $o_1,\dots,o_m$ to fix standardisation constants and the prior
       scale $g$; \textbf{discard} them
\State $\mathrm{status}(i,j) \gets \textsc{Undecided}$ for all pairs;
       $\mathrm{arrow}(i,j) \gets \textsc{false}$ for all ordered pairs
\State $\bar A(i,j) \gets 0$; $\bar E(i,j) \gets 0$
       \Comment{running maxima}
\For{each batch $B$}
  \State update the running Gram matrix with $B$
  \ForAll{pairs $(i,j)$ with $\mathrm{status}(i,j) = \textsc{Undecided}$}
    \State $A \gets \min_{S \in \Scal_k(i,j)} E^{(S)}(i,j)$
           \Comment{Schur complements of the Gram matrix}
    \State $\bar A(i,j) \gets \max\{\bar A(i,j), A\}$
    \If{$\bar A(i,j) \ge \lambda_A$}
      \State $\mathrm{status}(i,j) \gets \textsc{CertifiedEdge}$
      \State $Z_{ij} \gets \{\,\ell : \mathrm{status}(i,\ell) \text{ or }
             \mathrm{status}(j,\ell) = \textsc{CertifiedEdge} \,\}$
             \Comment{frozen}
      \State open direction processes for $(i,j)$ and $(j,i)$, starting from
             the next observation
    \EndIf
  \EndFor
  \ForAll{open ordered pairs $(i,j)$}
    \State update $\log E(i,j)$ by~\eqref{eq:direction} on the observations
           since the pair was opened
    \State $\bar E(i,j) \gets \max\{\bar E(i,j), E(i,j)\}$
    \If{$\bar E(i,j) \ge \lambda_D$} $\mathrm{arrow}(i,j) \gets \textsc{true}$
    \EndIf
  \EndFor
  \State \textbf{output} the three-valued graph; the analyst may stop here
\EndFor
\end{algorithmic}
\end{algorithm}

The output graph contains only certified edges. An edge carries an arrowhead
at an end where a direction certificate has fired and a circle otherwise, in
the mark notation of \citet{richardson2002ancestral}. \textbf{The absence of
an edge means undecided, not certified absent.}

\subsubsection{The whole-graph guarantee}\label{sec:theorem}

\begin{theorem}[Whole-graph anytime validity]\label{thm:main}
Suppose Assumptions~\ref{ass:markov} and~\ref{ass:sep} hold, the per-set
processes $E^{(S)}$ are e-processes for their point nulls, and---for the
direction certificates only---the pairs are generated by the linear
non-Gaussian model of Section~\ref{sec:direction} with the supremum
in~\eqref{eq:direction} attained. Then
\begin{equation}\label{eq:mainbound}
\Prob\left(
\begin{array}{c}
\exists\, t \ge 0 : \text{some edge certified by time $t$ is absent} \\[-1pt]
\text{from } \Gstar, \text{ or some arrowhead certified by time $t$
  is wrong}
\end{array}
\right) \;\le\; \alpha .
\end{equation}
In particular, for \emph{any} stopping time $\tau$---including one chosen by
inspecting the output---the graph reported at time $\tau$ satisfies
$\Prob(\text{it contains a false certified edge or arrowhead}) \le \alpha$.
\end{theorem}

\begin{proof}
Index the adjacency nulls by pairs and the direction nulls by ordered pairs;
both families are determined by $d$ and $k$ and are fixed before any
observation. Consider a single adjacency null for a pair $(i,j)$ that is
non-adjacent in $\Gstar$. By Lemma~\ref{lem:certifying}, $P \in
H^{\mathrm{sep}}_{ij}$; by Proposition~\ref{prop:adjacency}, $(A_t(i,j))$ is
an e-process for that null; so by~\eqref{eq:ville} the probability that
$A_t(i,j)$ ever reaches $\lambda_A = \binom{d}{2}/\alpha_A$ is at most
$\alpha_A / \binom{d}{2}$. Since the algorithm certifies the edge exactly on
that event, the probability that this pair is ever falsely certified is at
most $\alpha_A/\binom{d}{2}$.

Similarly, consider an ordered pair $(i,j)$ for which the arrowhead $i \to j$
is absent from $\Gstar$. Then $P \in H^{\mathrm{dir}}_{j \to i}$, which by
Proposition~\ref{prop:direction} (applied conditionally on $\Fcal_T$, see
Appendix~\ref{app:conditioning}) has $(E_t(i,j))$ as an e-process, and the
probability that it ever reaches $\lambda_D$ is at most
$\alpha_D / (d(d-1))$.

The event in~\eqref{eq:mainbound} is the union, over the at most
$\binom{d}{2}$ non-adjacent pairs and at most $d(d-1)$ absent arrowheads, of
these events. A union bound gives
$\binom{d}{2} \cdot \alpha_A/\binom{d}{2} + d(d-1) \cdot \alpha_D/(d(d-1))
= \alpha_A + \alpha_D = \alpha$.

The final statement follows because the event in~\eqref{eq:mainbound}
contains, for every stopping time $\tau$, the event that the graph at time
$\tau$ contains a false certified assertion; the bound is time-uniform, so no
separate argument per stopping time is needed.
\end{proof}

We stress the form of the quantifier. The theorem does not say ``for each
fixed $\tau$, the error at $\tau$ is at most $\alpha$'', which would be a
family of pointwise statements requiring a further union bound to combine.
It says that the probability of the trajectory \emph{ever} containing a false
certified assertion is at most $\alpha$, and every stopping-time statement is
a consequence. This is the difference between a fixed-sample guarantee applied
repeatedly and a time-uniform one, and Section~\ref{sec:exp-calibration}
measures how large that difference is in practice.

\begin{remark}[What the theorem does not cover]
Theorem~\ref{thm:main} bounds the probability of a false \emph{assertion}. It
says nothing about the pairs left undecided, and it does not bound the
probability that a true edge is missed---there is no such bound, and there
cannot be one without a faithfulness-type assumption. It also does not cover
the analyst's use of the output for anything else: an intervention chosen on
the basis of a certified arrowhead inherits the $\alpha$, but a downstream
effect estimate computed from the certified graph does not automatically
inherit anything, and would need its own analysis; the estimand there is the intervention-effect
problem of \citet{maathuis2009estimating}, and combining a certified graph
with a valid effect interval is a genuinely open question.
\end{remark}

\subsubsection{Complexity}\label{sec:complexity}

Per batch, the Gram-matrix update is $O(d^2)$ per observation. Recomputing
$A_t(i,j)$ for a pair costs $O(|\Scal_k| \cdot k^3)$, dominated by the Schur
complements, and $|\Scal_k| = \sum_{r \le k}\binom{d-2}{r} = O(d^k)$. Summed
over pairs, an update of the whole adjacency layer is $O(d^{k+2} k^3)$, with
no dependence on $n$: all the data enter through the sufficient statistic.
Only undecided pairs are updated, so the cost falls as the run proceeds.

The direction layer is the expensive part. Each open ordered pair requires a
generalised-Gaussian fit at every batch, which is an IRLS loop plus a scalar
optimisation over $\kappa$, and it does not reduce to a sufficient statistic
because the fit is nonlinear in the residuals. The cost is $O(d^2 \cdot
n_{\text{batch}} \cdot |Z| \cdot \text{iterations})$ per batch, and it is the
reason our experiments stop at $d = 11$. This is an implementation limit, not
a limit of the construction: the numerator's parameters may be refitted at any
frequency (refitting less often only costs power, since the numerator needs to
be $\Fcal_{s-1}$-measurable and nothing more), and the denominator's maximised
likelihood admits warm starts across batches. We have not pursued either
optimisation and regard scaling as unproven; see Section~\ref{sec:discussion}.

\subsubsection{On the multiplicity layer}\label{sec:multiplicity}

Theorem~\ref{thm:main} uses a union bound over a fixed family, and a union
bound is the least powerful admissible way to combine evidence. It is
reasonable to ask why we do not use a sharper procedure. E-value-based
family-wise error control has advanced considerably: \citet{hartog2026fwer}
show that e-closed testing with weighted arithmetic-mean local tests strictly
dominates e-Bonferroni, and their e-Holm procedure achieves this at $O(m \log
m)$ cost. Crucially, the arithmetic-mean local test is valid under
\emph{arbitrary} dependence between the base e-values, so the fact that all
our per-hypothesis processes are functions of the same running Gram matrix is
\emph{not} an obstacle. An earlier draft of this paper claimed it was; that
claim was wrong, and we correct it here.

The actual obstruction is different and, we think, more interesting. Our
certificates fire on the \emph{running maximum} $\bar A_t = \sup_{s \le t}
A_s$ of an e-process, because that is what a monitored run rejects on. A
running maximum is not an e-value. Ville's inequality bounds
$\Prob(\bar A_\infty \ge 1/\alpha)$ by $\alpha$, which makes $\bar A_t$ what
\citet{hartog2026fwer} call a \emph{pseudo} e-value, but $\Exp[\bar A_\tau]$
generally exceeds one, and closed testing takes e-values as inputs, not
pseudo e-values. Feeding running maxima into an e-closed procedure is not
licensed by the theory. \citet{hartog2026fwer} address this case directly and
recover a bound of $\alpha + O(\alpha^2 \log \alpha^{-1})$ for pseudo
e-values, but only under \emph{independence} across hypotheses---and that is
where our shared sufficient statistic does bite.

So the situation is: dependence is fine for e-closed testing with e-values;
running maxima are fine under independence; we have running maxima under
dependence, and no result covers the combination. The union bound over Ville
events is what remains available, and it is exact for what it does. Closing
this gap---a family-wise procedure for dependent pseudo e-values that improves
on Bonferroni---would immediately improve the power of every method in this
paper, and we regard it as the most valuable piece of the surrounding theory
that is missing.

Two weaker improvements are available now and we note them. First, the budget
split need not be uniform: any fixed weights $w_h > 0$ with $\sum_h w_h = 1$
give thresholds $1/(\alpha w_h)$ and the same bound, so prior knowledge about
which pairs are likely adjacent can be spent as unequal weights, decided
before the data arrive. Second, if the analyst wants false-discovery-rate
rather than family-wise control, the e-BH procedure of
\citet{wang2022false} applies to e-values under arbitrary dependence, as do
the martingale-based merging rules of \citet{vovk2024merging}, and would
replace the union bound directly---but only at a fixed stopping time, since it
again takes e-values rather than running maxima as inputs.

\subsubsection{An optional fourth verdict: certified negligible}
\label{sec:negligible}

Section~\ref{sec:principle} argued that there is no certifying null for
non-adjacency. There is, however, one for a weaker proposition, and analysts
who need it can have it at the cost of a stated margin. Fix $\delta > 0$ and
consider the interval null
\begin{equation}\label{eq:ropenull}
H^{\mathrm{big}}_{ij,S} \;:\quad |\rho_{ij \cdot S}| \ge \delta ,
\end{equation}
where $\rho_{ij\cdot S}$ is the partial correlation. Rejecting it for
\emph{every} $S$ in a chosen set certifies that the pair's association is
smaller than $\delta$ given each of those conditioning sets---a statement
about effect size, not about the graph. The confidence
sequence~\eqref{eq:cs} supplies the test directly: reject when the interval
excludes $[-\delta, \delta]$. The resulting verdict, \emph{certified
negligible at margin $\delta$}, is honest in a way that a deleted edge is not,
because $\delta$ is reported alongside it and the reader can judge whether an
association below $\delta$ matters for their purpose. It is also strictly less
than the claim a delete-on-non-rejection method makes, which is presumably why
that claim is popular. We implement this verdict but do not use it in the
experiments below, since the comparisons there are with methods that assert
absence outright.

\subsection{Latent confounders}\label{sec:latent}

The adjacency certificate does not require causal sufficiency; nothing in
Lemma~\ref{lem:certifying} assumes that all common causes are observed.
What changes under latent confounding is the interpretation of a certified
edge and the availability of the direction certificate.

If $\Gstar$ is a DAG over observed and latent variables and we observe only
$V$, the relevant object is the maximal ancestral graph over $V$
\citep{richardson2002ancestral}, and the relevant separation notion is
m-separation. Assumption~\ref{ass:sep} becomes the requirement that every
non-adjacent pair in the MAG has an m-separating set of size at most $k$ among
the observed variables---which, unlike the DAG case, is not guaranteed by
bounding an in-degree, and is why FCI searches Possible-D-SEP sets rather than
neighbourhoods. With that caveat, Lemma~\ref{lem:certifying} and
Proposition~\ref{prop:adjacency} go through verbatim with d-separation
replaced by m-separation, and a certified edge means ``$i$ and $j$ are
adjacent in the MAG'', that is, neither is separable from the other by any
observed set---which is compatible with a direct effect, a latent common
cause, or both.

We therefore implement CERT-FCI: run the adjacency certificates to obtain a
certified skeleton, apply the FCI orientation rules R1--R4 \citep[originally
][in the formulation of \citealp{zhang2008completeness}]{spirtes1995causal} to the certified edges only, and
report the resulting partial ancestral graph. Orientations produced by the FCI
rules inherit the adjacency certificates' guarantee only in the weak sense
that they are sound given a correct skeleton and correct sepsets; we do not
claim a family-wise bound for them, and we flag this as the principal gap in
the latent-variable extension. The direction certificate of
Section~\ref{sec:direction} is unavailable here, because
Proposition~\ref{prop:lingam} assumes no unobserved common cause of the pair.
Section~\ref{sec:exp-latent} reports what CERT-FCI recovers, and
Section~\ref{sec:exp-violations} reports what latent confounding does to the
direction certificate when it is applied anyway.
